\begin{textsample}{POS Dim 1 – vem\_ed\_33 – Score 19.00 – vem\_ed\_33/t175.txt}  \label{ex:f1_pos_076}
Boas práticas de Internacionalização em Casa Regiane Camargo, professora da Fatec Guaratinguetá e coordenadora de Projetos Colaborativos Internacionais em espanhol na AIES / Depes / CGESG, participou de mesa-redonda virtual com as colegas da Fatec Guaratinguetá: Patrícia Januária Barbosa, Tálita Guarino e Valdnéa Auxiliadora de Paula Santos.
A atividade, no dia 3 de dezembro de 2025, fez \textbf{parte} do VIII Seminário de Boas Práticas de Ensino e Aprendizagem (SBPEA) da Escola de Engenharia de Lorena da Universidade de São Paulo (EEL-USP).
Intitulada “Internacionalização em \textbf{casa}: relatos e boas práticas de aprendizagem \textbf{colaborativa} internacional – COIL”, a mesa-redonda abordou os PCIs realizados nas Fatecs, \textbf{seguindo} a abordagem mundialmente conhecida como COIL.
A gravação do evento encontra-se em VIII Seminário de Boas Práticas de Ensino e Aprendizagem (SBPEA) da EEL-USP.
Regiane fez uma introdução sobre os fundamentos dos PCIs / COIL e lembrou \textbf{que} podem \textbf{ser} realizados de maneira intra, inter ou transdisciplinar (por exemplo, envolvendo os Objetivos de Desenvolvimento \textbf{Sustentável}).
Destacou as etapas dos PCIs / COIL: quebra-gelo, \textbf{momento} cultural, \textbf{tarefa} principal, apresentação dos “entregáveis” e avaliação.
Tálita \textbf{compartilhou} vídeos com depoimentos de alunos envolvidos em PCIs orientados por ela, \textbf{destacando} a desenvoltura conquistada no \textbf{idioma} inglês e a aprendizagem sobre outras culturas.
Patrícia Januária Barbosa relatou o PCI realizado entre estudantes das Fatecs Bauru e Guaratinguetá e da SUNY Delhi (EUA). “Por meio das \textbf{interações} \textbf{síncronas}, nossos alunos perceberam \textbf{que} os colegas nos Estados Unidos atravessam desafios sociais semelhantes”.
Valdnéa abordou PCIs nas áreas de Negócios Internacionais e Comércio Exterior realizados em inglês e espanhol com Israel e Chile, respectivamente.
Os alunos israelenses do Sapir College \textbf{compartilharam} vídeo sobre seu cotidiano, mostrando os abrigos antibombas.
A faculdade fica a apenas 4 km da fronteira com a Faixa de Gaza, e os alunos só vão presencialmente 4 vezes por semestre à instituição. “\textbf{Quanto} mais conhecemos sobre nosso parceiro, melhor a comunicação e a negociação”.
O PCI com o Chile, por sua vez, teve o objetivo de proporcionar uma \textbf{compreensão} \textbf{prática} do \textbf{processo} de exportação entre Brasil e Chile.

% matched lemmas: casa, colaborativo, compartilhar, compreensão, destacar, idioma, interação, momento, parte, processo, prático, quanto, que, seguir, ser, sustentável, síncrono, tarefa
\end{textsample}
